% Generated draft: requires mathematical and source review.
\documentclass[11pt]{article}
\usepackage[margin=25mm]{geometry}
\usepackage{fontspec}
\setmainfont{Latin Modern Roman}
\usepackage{amsmath,amssymb,mathtools,unicode-math}
\setmathfont{Latin Modern Math}
\usepackage{enumitem,needspace,xurl,hyperref}
\setlength{\parindent}{0pt}
\setlength{\parskip}{5pt}
\setlength{\emergencystretch}{3em}
\newcommand{\N}{\mathbb{N}}
\newcommand{\Z}{\mathbb{Z}}
\newcommand{\Q}{\mathbb{Q}}
\newcommand{\R}{\mathbb{R}}
\newcommand{\C}{\mathbb{C}}
\newcommand{\F}{\mathbb{F}}
\newcommand{\A}{\mathbb{A}}
\newcommand{\PP}{\mathbb{P}}
\newcommand{\E}{\mathbb{E}}
\newcommand{\eps}{\varepsilon}
\newcommand{\dd}{\,\mathrm{d}}
\begin{document}
\section*{UnsolvedMath ID: 30003258}
\textbf{Source problem code:} catalogue entry
\subsection{Definitions and mathematical statement}
The supplied record is a composite of several questions, with statement status recorded as \texttt{unrecoverable}. Its supplied verification note says that unrelated questions were concatenated and that no clean statement was promoted. The following separates the recoverable questions and states the remaining qualifications explicitly.

\begin{enumerate}
\item Let \(\C(t)\) be the rational function field in one indeterminate over \(\C\), and let
\[
\mathcal O_t=\C[t]_{(t)}=\left\{\frac{a(t)}{b(t)}:a,b\in\C[t],\ b(0)\ne0\right\}.
\]
This is the valuation ring of the valuation given by order of vanishing at \(t=0\). For a field \(K\), a subset \(S\subseteq K\) is diophantine with parameters from \(K\) if there are integers \(m\ge0\), \(r\ge1\), and polynomials \(P_1,\ldots,P_r\in K[X,Y_1,\ldots,Y_m]\) such that
\[
x\in S\quad\Longleftrightarrow\quad
\exists\mathbf y\in K^m\ \bigwedge_{i=1}^{r}P_i(x,\mathbf y)=0.
\]
The first question asks whether \(\mathcal O_t\) is diophantine in \(\C(t)\). The record does not specify permitted parameters or a coefficient subring for the defining equations; the definition above describes the version allowing arbitrary field parameters, without asserting that this is the intended convention.

\item For \(K=\C(s,t)\) or \(K=\F_p(s,t)\), where \(s,t\) are algebraically independent and \(p\) is prime, the source asks whether the condition
\[
\left.\left.f\right|_{s=0}\right|_{t=0}=0
\]
is diophantine in \(K\). The displayed expression involves successive specializations, which are not defined for every rational function. One possible precise interpretation, for \(k=\C\) or \(k=\F_p\), uses
\[
A=k(t)[s]_{(s)},\qquad \rho_s:A\longrightarrow k(t),\qquad \rho_s(g)=g|_{s=0},
\]
and the subset
\[
Z_k=\left\{f\in A:\rho_s(f)\in k[t]_{(t)}\ \text{and}\ \rho_s(f)|_{t=0}=0\right\}\subseteq k(s,t).
\]
Under this interpretation the question would be whether \(Z_k\) is diophantine in \(k(s,t)\). This is only a conditional reconstruction: the supplied record does not establish the intended specialization domain or whether the target is this zero set, a valuation ring, or a relation encoding the place. Permitted parameters and the quantification over primes \(p\) are also unspecified.

\item The record asks whether an infinite finitely axiomatizable field exists. For a specified first-order language \(L\) and an \(L\)-structure \(K\), its complete theory \(\operatorname{Th}_L(K)\) is finitely axiomatizable if there is a finite set \(\Sigma\) of \(L\)-sentences such that, for every \(L\)-structure \(M\),
\[
M\models\Sigma\quad\Longleftrightarrow\quad M\models\operatorname{Th}_L(K).
\]
With this meaning and with the pure ring language \(L=\{0,1,+,-,\cdot\}\), the question is whether there exists an infinite field \(K\) for which \(\operatorname{Th}_L(K)\) is finitely axiomatizable. The supplied record does not explicitly specify the language or the intended meaning of finite axiomatizability, so this formulation remains conditional on that interpretation.

\item Write \(\Q^*\succeq\Q\) for an elementary extension in the intended language, and let \(v\) be a valuation on \(\Q^*\). A henselization \((\Q^*)^h_v\) is an algebraic henselian extension of the valued field with the usual minimality property, unique up to valued-field isomorphism over \(\Q^*\). A field is separably closed if it has no proper finite separable algebraic extension. The source asks whether \((\Q^*)^h_v\) must be separably closed for a valuation and elementary extension ``as in the report.'' The referenced hypotheses on \(v\), on the extension, and on the language are absent. Consequently, this cannot be stated as an unconditional assertion about every valuation on every elementary extension of \(\Q\).
\end{enumerate}

The supplied status is \emph{open}; this is a source label, not an independently verified mathematical conclusion. These questions form a statement-only draft requiring human mathematical review.
\subsection{Short English statement}
Permanent UnsolvedMath ID \(30003258\) combines questions about diophantine definitions of valuation-related sets, the existence of an infinite finitely axiomatizable field, and separable closedness of a henselization of an elementary extension of \(\Q\). Some definitions can be made explicit, but the specialization target, parameter conventions, axiomatization language, and hypotheses for the henselization question remain unresolved.
\subsection{Sources}
\begin{itemize}
\item Supplied local UnsolvedMath registry record, permanent ID \(30003258\), title \emph{Composite of Unrelated Field-Theory Questions}, exact source problem code \texttt{catalogue entry}. The project registry is \texttt{data/unsolvedmath/problems.json}, with additional existing records in \texttt{catalog/problems.json}. Attribution is to the Oberwolfach Reports workshop open-problem collection; no individual proposer is supplied. The proposed year is recorded as 2016.
\item Supplied source citation: \emph{Definability and Decidability Problems in Number Theory} (2017), Oberwolfach Reports, DOI 10.4171/owr/2016/49, pp. 2793--2866. \url{https://doi.org/10.4171/owr/2016/49}. This is the report citation supplied by the record; the relevant passages have not been independently consulted here.
\item Supplied literature source: \emph{Finite undecidability in NIP fields}. \url{https://www.cambridge.org/core/journals/journal-of-symbolic-logic/article/finite-undecidability-in-nip-fields/FD28F4C52FA9926B68BF4278DDEE2751}.
\item Supplied literature source: \emph{Oberwolfach Reports 2017}. \url{https://ems.press/content/serial-article-files/46652}.
\item Literature assessment dated exactly \texttt{2026-08-21}: the supplied curation labels the record open, describes it as combining several definability and finite-axiomatizability questions, and states that current literature still identifies the infinite finitely axiomatizable field problem as open. This is supplied dated curation, not an independently verified assessment. The source itself requests independent verification of time-sensitive status and cited claims.
\item The supplied statement-verification note reports that the official problem-session list reveals a concatenation of unrelated questions. It records statement status as \texttt{unrecoverable}, promotes no clean statement, and retains the original record publicly for provenance.
\end{itemize}
\end{document}
